% Pursuing Stacks (cotes 134-2 to 134-8): archive cote 134-5, site page 44 (PDF page 45), typescript page 231.
% typescript. Draft compared: ps4.tex:2144-2221.
% AI-transcribed by Claude Opus 5.5 (claude-opus-5-5, adaptive thinking, effort high), Message Batches API, batch ps2 (msgbatch_01NVwr1Yuo8qJf5rykeGBXPK), 2026-11-12; not yet checked by eye.
\scanpage{134-5/44}
\mspage{231}
a mihimum of technical work,  needed for keeping out of the uneasiness
of ``thin air conjecturing''\hand{)}. An attentive reader will have felt occasiona%
-ly this implicit presence of $\underline{U}$,  for instance in the definition of the
basic modelizer (Cat) (which is,  as all modelizers,  a ``large'' category)\hand{,}
and in our occasional reference to various categories as being ``small'' or
``large''.  He will have noticed that whereas modelizers are by necessity
large categories (justas (Hot) itself,  whose set of isomorphism classes
of objects is large),  test categories are supposed to be small (and of-%
ten even to be in (Cat)\hand{)} -  with the effect that $A^{\wedge}$ ,  the category of
presheaves on A, is automatically a $\underline{U}$-category  (which would not be the
case if A was merely assumed to be a $\underline{U}$-category).

An \uline{asphericity} \uline{structure} (with re\over{p}{\unsure{s}}pect to the \unsure{badic|basic} localizer $\underline{W}$ )
on the $\underline{U}$-category M consists of a subset
\[ \text{(1)} \qquad M_{as} \subset \text{Ob } M \quad , \]
whose elements will be called the \uline{aspheric objects} of M  (more specifical%
ly,  the $\underline{W}$-aspheric objects,  if confusion may arise),  this subset being
submitted to the following condition:

(Asstr)  There exists an object A in (Cat),  and a functor $i: A \longrightarrow M$ ,
satisfying the following two conditions:
(i)  For any a in A,  $i(a) \in M_{as}$ ,  i.e.  i factors through the
full subcategory (also denoted by $M_{as}$) of M defined by $M_{as}$.
(ii)  Let
\[ i^{*}:\ M \longrightarrow A^{\wedge} \quad , \quad x \mapsto i^{*}(x) = (a \mapsto \text{Hom}(i(a),x)) \ , \]
\inserted[typed]{then we \xout{xatx}} \xout{xxxxxxxxxxxx}  have
\[ \text{(2)} \qquad M_{as} = (i^{*})^{-1}(A^{\wedge}_{as}) \quad , \]
where $A^{\wedge}_{as}$ is the subset of Ob $A^{\wedge}$ of all ($\underline{W}$-)aspheric objects
of $A^{\wedge}$,  i.e. the \xout{objects} presheaves F on A such that the
object $A_{/F}$  of (Cat) is $\underline{W}$-aspheric.

In other words, for x in M,  we have the equivalence
\[ \text{(2 bis)} \qquad x \in M_{as} \Longleftrightarrow i^{*}(x) \in A^{\wedge}_{as} \ \text{ i.e. } \ A_{/x}\ (\overset{def}{=} A_{/i^{*}(x)}) \in (\text{Cat})_{as} , \]
where $(\text{Cat})_{as}$ is the subset of Ob (Cat) made up with all $\underline{W}$-aspheric ob-%
jects of (Cat)  -  i.e. objects X such that  $X \hand{$\to \Delta_{\bullet}$}$ \hand{(}= final object\hand{)}
be in $\underline{W}$ .

Thus an asphericity structure on M can always be defined by a func-%
tor  \xout{ixxix}
\[ i:\ A \longrightarrow M \quad , \]
with A in (Cat),  by the formula  (2);  and conversely,  any such functor

\margin{44}
